% C015 椭圆焦点三角形四心坐标性质——直观解释
\begin{explanationbox}
\textbf{一、为什么这道结论值得记？}

先忘记椭圆：任何三角形，只要底边两个顶点是
$(-c,0)$ 和 $(c,0)$，第三个顶点是 $P(x_0,y_0)$，
就天然具备\textbf{横向对称、底边水平}这两个便利条件。
椭圆又额外送来一个条件：$PF_1+PF_2=2a$ 恒定。
因此，一般三角形的四心问题，在这里都能被大幅简化。

\textbf{二、把四个中心分别想成四种动作}
\begin{enumerate}
 \item \textbf{重心：三个点求平均。} 两焦点横坐标一正一负，直接抵消；
 因此 $G$ 就是把 $P$ 的坐标各缩小为原来的 $1/3$。
 从原点 $O$ 看，$O,G,P$ 共线，且 $OG:OP=1:3$。

 \item \textbf{内心：先找底边的切点，再找高度。}
 想象内切圆压在水平的底边 $F_1F_2$ 上，切点为 $T$。
 圆心 $I$ 一定在 $T$ 的正上方或正下方，故 $x_I=x_T$。
 从三角形同一顶点向圆作的两段切线等长，
 可以得到 $F_1T-F_2T=PF_1-PF_2$；而
 $F_1T-F_2T=2x_T$，所以
 \[
 x_I=\frac{PF_1-PF_2}{2}.
 \]
 椭圆焦半径之差恰为 $2ex_0$，故 $x_I=ex_0$。
 又因为三角形面积既等于 $c|y_0|$，也等于半周长乘内切圆半径，
 即 $c|y_0|=(a+c)r$，所以
 \[
 |y_I|=r=\frac{c|y_0|}{a+c}.
 \]
 这两步就得到了整个内心，\textbf{没有必要死记内心的加权坐标公式}。

 \item \textbf{垂心：沿着 $P$ 的竖线走。}
 底边在 $x$ 轴上，过 $P$ 的高一定垂直于 $x$ 轴，所以垂心与 $P$
 的横坐标天然相同：$x_H=x_0$。
 接下来只需再找一条高线来确定它的纵坐标。

 \item \textbf{外心：沿着 $y$ 轴找。}
 外心到 $F_1,F_2$ 距离相等，必须在 $F_1F_2$ 的垂直平分线——$y$ 轴上。
 \textbf{这与三角形是不是等腰毫无关系。}
 设 $O'=(0,t)$，只需要令 $O'P=O'F_2$，就能求 $t$。
\end{enumerate}

\textbf{三、为何椭圆让内心特别简单？}

对任意一个以上述两点为底边的三角形，都有
\[
 x_I=\frac{PF_1-PF_2}{2},\qquad
 y_I=\frac{2cy_0}{PF_1+PF_2+2c}.
\]
椭圆第一定义使 $PF_1+PF_2=2a$，
而\textbf{距离平方作差}又给出 $PF_1-PF_2=2cx_0/a$。
于是内心就成为对 $P$ 的固定比例缩放：
\[
 (x_0,y_0)\longmapsto
 \left(\frac ca x_0,\frac{c}{a+c}y_0\right).
\]
这是一种\textbf{横、纵缩放比例不同}的线性变换，
不是一般的等比例缩小。因此内心轨迹还是椭圆，但形状和原椭圆不同。

\textbf{四、两个特别好记的几何现象}
\begin{enumerate}
 \item \textbf{垂心随 $P$ 在同一竖线上。}
 当 $x_0=c$ 时，$PF_2$ 与底边垂直，三角形在 $F_2$ 处成直角，
 因此垂心就是 $F_2$；代入公式也正好得 $H=(c,0)$。
 当 $x_0=-c$ 时同理，$H=F_1$。

 \item \textbf{外心可能位于底边的另一侧。}
 令 $O'=(0,t)$，则
 \[
 2y_0t=x_0^2+y_0^2-c^2=OP^2-c^2.
 \]
 若 $OP=c$，以 $F_1F_2$ 为直径的圆恰好通过 $P$；
 由直径所对的圆周角为直角可知 $\angle F_1PF_2=90^\circ$，
 此时外心恰为椭圆中心 $O$。
 若 $OP<c$，顶角 $\angle F_1PF_2$ 是钝角，外心落在与 $P$ 相反的一侧。
\end{enumerate}

\textbf{五、两个中心之间还有一条最重要的捷径}

一般三角形有\textbf{欧拉线性质}：重心 $G$ 在外心 $O'$ 与垂心 $H$ 的连线上，
且 $\overrightarrow{O'G}=\tfrac13\overrightarrow{O'H}$。
本题因为 $G=P/3$，立即得到
\[
 H=3G-2O'=P-2O'.
\]
换句话说：\textbf{先求外心，再把外心纵坐标乘 $2$ 从 $P$ 的纵坐标中减去，垂心立刻出现。}
这比重新联立两条高线更省步骤。

\textbf{六、与圆的比较（注意不能直接令 $a=b$）}

任何圆的内切圆切线长定理、三角形重心、垂心、外心、欧拉线公式在这里都可以用；
区别在于椭圆又提供 $PF_1+PF_2=2a$，让内心变成一个固定的坐标变换。
\textbf{不能直接取 $a=b$ 来得到所谓“圆的焦点三角形四心”：}
此时 $c=0$、两个焦点重合，底边长度为零，三角形已经退化。

\textbf{七、图形的极限变化}

当 $P$ 趋近长轴端点时 $y_0\to0$，三角形越来越扁：
内切圆半径 $r\to0$，而垂心与外心的纵坐标绝对值趋于无穷大。
这不是公式出错，而是接近共线时高线及中垂线交点对微小形变非常敏感。
\end{explanationbox}
